\documentclass[11pt,a4paper]{article}

% -----------------------------------------------------
% Packages
% -----------------------------------------------------
\usepackage[margin=1in]{geometry}
\usepackage{hyperref}
\usepackage{enumitem}
\usepackage{graphicx}
\usepackage{xcolor}
\usepackage{titling}

% -----------------------------------------------------
% bvraghav-tiet-ucs503p-article.sty
% -----------------------------------------------------

% Variable: Lab Instructor
\makeatletter
\newcommand{\BvrSetLabInstructor}[1]{%
  \gdef\Bvr@LabInstructor{#1}%
}

% Default
\newcommand{\Bvr@LabInstructor}{Dr. Raghav B. Venkataramaiyer}

% Public accessor
\newcommand{\BvrLabInstructorValue}{\Bvr@LabInstructor}
\makeatother

% -----------------------------------------------------
% Title formatting
% -----------------------------------------------------

\pretitle{%
  \begin{center}
  \bfseries
  \fontsize{18bp}{18bp}\selectfont
}

\makeatletter
\newcommand{\BvrSubtitle}[1]{%
  \posttitle{%
    \par
    \end{center}
    \begin{center}
    \large #1
    \end{center}
    \vskip 0.5em
  }%
}
\makeatother

% -----------------------------------------------------
% Section hints
% -----------------------------------------------------

\newcommand{\BvrHint}[1]{%
  {\normalfont\normalsize\itshape\hfill #1}%
}

% -----------------------------------------------------
% Project Title
% -----------------------------------------------------

\title{
  \textbf{SAGE} ---
  \textit{Smart Adaptive Gating Engine for Cost-Efficient Code-Model Selection}
}

% -----------------------------------------------------
% Lab Instructor
% -----------------------------------------------------

\BvrSetLabInstructor{Ms. Paramveer Sidhu}

% -----------------------------------------------------
% Subtitle
% -----------------------------------------------------

\BvrSubtitle{%
  Project Proposal for UCS503P \\
  Submitted to: \BvrLabInstructorValue \\[2pt]
  \bfseries Thapar Institute of Engineering and Technology%
}

% -----------------------------------------------------
% Student Details
% -----------------------------------------------------

\author{% 
  \begin{tabular}{c@{\qquad}c@{\qquad}c}
    \textbf{Krishna Gautam} & \textbf{Prateek Yadav} & \textbf{Harshvardhan} \\[3pt]
    Roll No: \texttt{1024030733} & Roll No: \texttt{1024031149} & Roll No: \texttt{1024030744} \\
    \texttt{kgautam\_be24@thapar.edu} & \texttt{pyadav\_be24@thapar.edu} & \texttt{hharshvardhan\_be24@thapar.edu}
  \end{tabular}%
}

\date{\today}

% -----------------------------------------------------
% Document
% -----------------------------------------------------

\begin{document}

\maketitle

% -----------------------------------------------------
\section{Higher-order goal%
  \texorpdfstring{\BvrHint{\ldots secondary framing}}{}}
% -----------------------------------------------------

This project aims to improve the efficiency of AI-assisted
software development by automatically selecting an appropriate
code-generation model for each coding request. While modern
large language models can solve a wide range of programming
tasks, using the strongest model for every request can result
in unnecessary cost, latency, and computational usage. The
immediate engineering objective is to translate
\emph{cost-efficient and quality-aware model selection} into
a measurable and deployable software solution.

% -----------------------------------------------------
\section{Time-to-value%
  \texorpdfstring{\BvrHint{\ldots primary heuristic}}{}}
% -----------------------------------------------------

\begin{itemize}[leftmargin=0.7cm]

    \item We will deliver meaningful increments quickly by
    first implementing a baseline router that classifies
    coding prompts into difficulty levels and maps them to
    suitable Claude models.

    \item We will prioritize measurable experimentation
    through short iterations, comparing the routing system
    against a baseline that always uses a stronger model.

    \item Success is defined by measurable improvement in
    cost and latency while maintaining acceptable coding
    quality and correctness.

\end{itemize}

% -----------------------------------------------------
\section{Problem Statement}
% -----------------------------------------------------

AI coding assistants increasingly rely on powerful large
language models to generate, debug, explain, and refactor
software. However, not every coding request requires the
same level of reasoning capability. A simple syntax question
and a complex multi-file debugging task may both be sent to
the same expensive model even though their computational
requirements are very different. Common pain points include:

\begin{itemize}[leftmargin=0.7cm]

    \item \textbf{Over-provisioning:} simple coding tasks may
    consume the capabilities and cost of a high-end reasoning
    model unnecessarily.

    \item \textbf{Cost inefficiency:} repeatedly using a
    stronger model increases token consumption and overall
    inference cost for workloads containing many easy or
    moderate tasks.

    \item \textbf{Latency:} stronger models may introduce
    additional response time when a faster model would have
    been sufficient.

    \item \textbf{Difficult model selection:} users and
    applications generally have to decide manually which
    model should handle a particular coding request.

    \item \textbf{Quality-risk tradeoff:} selecting a cheaper
    model can reduce cost, but selecting one that is too weak
    for a difficult task can reduce correctness and usefulness.

\end{itemize}

\textbf{Why this matters now (contextual relevance):}
As AI-assisted programming becomes a regular part of software
development, efficient model utilization becomes increasingly
important. A routing layer that can automatically match task
complexity with model capability can reduce unnecessary
inference cost without simply sacrificing output quality.

% -----------------------------------------------------
\section{Proposed Solution%
  \texorpdfstring{\BvrHint{\ldots Software Proposition}}{}}
% -----------------------------------------------------

\subsection{Overview}

Build a software-based ``gating engine'' that analyzes a
coding request before sending it to a Claude model. The system
will estimate the complexity of the task, select an appropriate
model tier, forward the request, and record the result for
evaluation.

\begin{itemize}[leftmargin=0.7cm]

    \item \textbf{Coding prompt intake:} accept a programming
    request together with optional context such as language,
    code snippets, files, and task type.

    \item \textbf{Prompt analysis:} extract signals related
    to coding difficulty, reasoning requirements, context
    size, debugging complexity, and expected output.

    \item \textbf{Gating engine:} assign the request to an
    appropriate difficulty/model tier.

    \item \textbf{Model execution:} send the request to the
    selected Claude model through the appropriate API.

    \item \textbf{Evaluation and logging:} record routing
    decisions, token usage, latency, and task outcomes for
    comparison against baseline strategies.

\end{itemize}

% -----------------------------------------------------
\subsection{Core workflow}
% -----------------------------------------------------

\begin{enumerate}[leftmargin=0.7cm]

    \item User submits a coding prompt to the SAGE gateway.

    \item The system extracts relevant task characteristics,
    such as task type, code/context size, reasoning depth,
    and expected implementation complexity.

    \item The gating layer assigns a difficulty score or
    difficulty class such as Easy, Moderate, Hard, or
    Very Hard.

    \item The routing policy maps the difficulty class and
    confidence to a suitable Claude model.

    \item The selected model generates the response.

    \item The system records cost-related and quality-related
    metrics and can identify cases where a stronger model would
    have been more appropriate.

\end{enumerate}

% -----------------------------------------------------
\subsection{Operational constraints for an educational setting}
% -----------------------------------------------------

\begin{itemize}[leftmargin=0.7cm]

    \item The initial implementation will focus on coding
    prompts rather than attempting to route every possible
    natural-language task.

    \item The system will use commercially available model APIs
    and will not attempt to train a large language model from
    scratch.

    \item Routing decisions should be lightweight so that the
    cost and latency introduced by the gating layer do not
    outweigh the savings achieved through model selection.

    \item The project will emphasize measurable evaluation,
    reproducibility, and system correctness rather than
    claiming a new state-of-the-art language model.

\end{itemize}

% -----------------------------------------------------
\section{Solution Approach%
  \texorpdfstring{\BvrHint{\ldots Engineering focus}}{}}
% -----------------------------------------------------

This is an engineering-focused AI systems project. The main
focus is reliable model routing, measurable cost-quality
optimization, and evaluation rather than training a new
foundation model.

% -----------------------------------------------------
\subsection{Prompt difficulty classification%
  \texorpdfstring{\BvrHint{\ldots routing signal}}{}}
% -----------------------------------------------------

The gating engine will classify coding requests using a
combination of observable task characteristics:

\begin{itemize}[leftmargin=0.7cm]

    \item \textbf{Task type:} explanation, code generation,
    debugging, optimization, refactoring, testing, or
    architecture.

    \item \textbf{Reasoning requirement:} whether the task
    requires straightforward generation or multi-step
    reasoning and problem solving.

    \item \textbf{Context complexity:} approximate amount and
    structure of code or repository context supplied to the
    model.

    \item \textbf{Implementation complexity:} number of
    components, dependencies, edge cases, or interacting
    requirements.

    \item \textbf{Routing confidence:} confidence in the
    selected difficulty tier, allowing uncertain requests
    to be routed more conservatively.

\end{itemize}

% -----------------------------------------------------
\subsection{Gating and model-selection policy}
% -----------------------------------------------------

The first version will use a transparent routing policy that
maps difficulty tiers to model tiers. Easy requests will be
directed toward a faster and lower-cost model, moderate
requests toward a balanced model, and difficult requests
toward a stronger reasoning-capable model.

The policy will be evaluated rather than assuming that the
strongest model is always optimal. If the system detects low
confidence or a task that exceeds the expected capability of
the selected model, a fallback mechanism can escalate the
request to a stronger model.

% -----------------------------------------------------
\subsection{Evaluation architecture%
  \texorpdfstring{\BvrHint{\ldots quality-aware routing}}{}}
% -----------------------------------------------------

A separate evaluation layer will compare routing decisions
against baseline strategies. The evaluation layer will
consider:

\begin{itemize}[leftmargin=0.7cm]

    \item correctness of generated code,
    \item test-case pass rate where executable tests are
    available,
    \item token consumption and estimated inference cost,
    \item response latency,
    \item routing accuracy,
    \item and the percentage of tasks successfully handled
    without escalation.

\end{itemize}

This allows the project to measure whether lower-cost routing
actually preserves useful coding performance.

% -----------------------------------------------------
\subsection{System architecture%
  \texorpdfstring{\BvrHint{\ldots gateway-based solution}}{}}
% -----------------------------------------------------

A typical architecture will contain the following components:

\begin{itemize}[leftmargin=0.7cm]

    \item \textbf{Client/API layer:} accepts coding prompts
    and optional code or repository context.

    \item \textbf{Prompt analyzer:} extracts task and
    complexity signals.

    \item \textbf{Gating engine:} calculates the difficulty
    tier and routing confidence.

    \item \textbf{Model router:} maps the decision to the
    selected Claude model and forwards the request.

    \item \textbf{Evaluation and logging layer:} stores
    routing decisions, latency, usage, and evaluation results.

    \item \textbf{Dashboard or reporting layer:} presents
    model utilization, cost savings, quality, and routing
    performance.

\end{itemize}

% -----------------------------------------------------
\subsection{CI/CD and fast delivery enabler}
% -----------------------------------------------------

CI/CD will be used to:

\begin{itemize}[leftmargin=0.7cm]

    \item Automatically build and run unit and integration
    tests on every change.

    \item Validate routing rules and API integrations before
    deployment.

    \item Produce repeatable deployment artifacts for a
    staging environment.

    \item Enable safe and frequent improvements to the
    routing policy and evaluation pipeline.

\end{itemize}

% -----------------------------------------------------
\section{Evaluation Criterion%
  \texorpdfstring{\BvrHint{\ldots measurable and attributable}}{}}
% -----------------------------------------------------

We define success using measurable metrics tied directly to
model-selection efficiency and coding quality.

% -----------------------------------------------------
\subsection{Primary evaluation metric}
% -----------------------------------------------------

\textbf{Cost-quality efficiency:} the primary objective is to
reduce average model inference cost while maintaining coding
quality close to a strong-model baseline.

\begin{itemize}[leftmargin=0.7cm]

    \item \textbf{Target:} demonstrate a measurable reduction
    in average inference cost compared with always routing
    requests to the strongest available model, without a
    significant decrease in correctness.

    \item \textbf{Attribution:} compare identical benchmark
    prompts under the baseline and SAGE routing policies and
    record model usage, token consumption, latency, and
    evaluation scores.

\end{itemize}

% -----------------------------------------------------
\subsection{Secondary evaluation metrics}
% -----------------------------------------------------

\begin{itemize}[leftmargin=0.7cm]

    \item \textbf{Code correctness:} percentage of generated
    solutions that pass predefined tests where executable
    test cases are available.

    \item \textbf{Quality retention:} compare SAGE responses
    against responses generated by a strong-model baseline
    using task-specific evaluation criteria.

    \item \textbf{Cost reduction:} percentage reduction in
    estimated model inference cost relative to the
    always-strong-model baseline.

    \item \textbf{Latency reduction:} average and median
    response-time improvement resulting from routing simpler
    tasks to faster models.

    \item \textbf{Routing accuracy:} percentage of tasks for
    which the selected model tier is judged sufficient for
    the task.

    \item \textbf{Escalation rate:} percentage of requests
    initially routed to a lower tier that require escalation
    to a stronger model.

\end{itemize}

% -----------------------------------------------------
\subsection{Benchmark validation plan%
  \texorpdfstring{\BvrHint{\ldots high-level}}{}}
% -----------------------------------------------------

\begin{itemize}[leftmargin=0.7cm]

    \item Build a representative benchmark of coding prompts
    covering explanation, generation, debugging, refactoring,
    algorithms, testing, and multi-step engineering tasks.

    \item Divide prompts into difficulty categories using
    clearly defined criteria and, where possible, human
    validation.

    \item Run the same benchmark using an always-strong-model
    baseline and the SAGE routing policy.

    \item Compare correctness, cost, latency, model
    utilization, and escalation behavior.

    \item Analyze failure cases where the router selected a
    model that was too weak or unnecessarily strong.

\end{itemize}

% -----------------------------------------------------
\section{Scalability%
  \texorpdfstring{\BvrHint{\ldots with foresight}}{}}
% -----------------------------------------------------

The proposition should scale in both theoretical and
practical senses.

\begin{enumerate}[leftmargin=0.7cm]

    \item \textbf{Scalable (theoretically):} The routing
    architecture should support increasing numbers of coding
    requests by:

    \begin{itemize}[leftmargin=0.7cm]
        \item keeping the routing layer lightweight,
        \item separating prompt analysis from model execution,
        \item using asynchronous request handling where
        appropriate,
        \item and storing evaluation data independently from
        the request path.
    \end{itemize}

    \item \textbf{Operationally deployable (practically):}
    The system should be deployable as a gateway service:

    \begin{itemize}[leftmargin=0.7cm]
        \item containerized backend,
        \item managed database or logging store,
        \item configurable model-provider credentials,
        \item and CI/CD pipeline for staging and production.
    \end{itemize}

\end{enumerate}

% -----------------------------------------------------
\section{Engine availability heuristic}
% -----------------------------------------------------

A mature set of software components and AI services is
available to implement the proposed system:

\begin{itemize}[leftmargin=0.7cm]

    \item Backend framework for exposing the routing API.

    \item Claude model APIs for model execution.

    \item Standard machine-learning or rule-based techniques
    for prompt classification and scoring.

    \item Database or structured storage for prompts,
    routing decisions, evaluation results, and usage data.

    \item Testing frameworks for unit, integration, and
    benchmark evaluation.

    \item Containerization and CI/CD tooling for repeatable
    deployment.

\end{itemize}

We will use these mature components to focus on routing
strategy, evaluation, correctness, and measurable
cost-quality optimization rather than building model
infrastructure from scratch.

% -----------------------------------------------------
\section{Project scope and deliverables%
  \texorpdfstring{\BvrHint{\ldots iteration-friendly}}{}}
% -----------------------------------------------------

\subsection{Initial deliverable%
  \texorpdfstring{\BvrHint{\ldots first iteration}}{}}

\begin{itemize}[leftmargin=0.7cm]

    \item Coding prompt gateway/API.

    \item Initial difficulty classification scheme with
    transparent routing rules.

    \item Integration with multiple Claude model tiers.

    \item Basic logging of selected model, latency, and
    token usage.

    \item A small benchmark dataset covering multiple coding
    difficulty levels.

    \item Baseline comparison against always using a strong
    model.

    \item CI: build, unit tests, linting, and routing tests.

\end{itemize}

% -----------------------------------------------------
\subsection{Subsequent deliverables}
% -----------------------------------------------------

\begin{itemize}[leftmargin=0.7cm]

    \item Improved difficulty classifier based on benchmark
    results.

    \item Confidence-aware routing and automatic escalation
    for uncertain or failed tasks.

    \item Expanded evaluation suite for code correctness,
    cost, latency, and quality retention.

    \item Dashboard/report showing model distribution,
    estimated cost savings, routing accuracy, and failure
    cases.

    \item Support for richer repository and multi-file
    context analysis.

    \item Optimization of the routing policy using observed
    evaluation results.

\end{itemize}

% -----------------------------------------------------
\section{Risks and mitigations}
% -----------------------------------------------------

\begin{itemize}[leftmargin=0.7cm]

    \item \textbf{Incorrect difficulty classification:}
    define explicit classification criteria, maintain a
    benchmark set, and use confidence-aware fallback to a
    stronger model.

    \item \textbf{Quality degradation:} evaluate routed
    responses against a strong-model baseline and prioritize
    correctness over cost when the router is uncertain.

    \item \textbf{Routing overhead:} keep the gating layer
    lightweight so that analysis cost and latency do not
    eliminate the benefits of model selection.

    \item \textbf{Changing model capabilities or pricing:}
    keep model-selection policies configurable rather than
    hard-coding assumptions about a particular model.

    \item \textbf{Evaluation ambiguity:} define benchmark
    tasks and objective metrics such as test pass rate,
    token usage, latency, and estimated cost before
    comparing routing strategies.

    \item \textbf{API availability and limits:} use a staging
    environment, implement error handling, and maintain
    clear separation between routing logic and provider
    integration.

\end{itemize}

% -----------------------------------------------------
\section{Summary}
% -----------------------------------------------------

This project proposes SAGE, a deployable gating engine for
cost-efficient selection of Claude models for coding tasks.
It addresses the practical problem of using unnecessarily
powerful models for simple requests while avoiding quality
loss on difficult programming tasks. The project meets the
course criteria by emphasizing rapid engineering iterations,
CI/CD, measurable evaluation, and scalability readiness.

The key contribution is not a new language model, but a
measurable model-routing system that attempts to achieve a
better tradeoff between coding quality, inference cost, and
response latency. By comparing SAGE against an
always-strong-model baseline, the project will demonstrate
whether intelligent model selection can provide meaningful
efficiency gains while preserving useful coding performance.

% -----------------------------------------------------
\end{document}
